\section{Conclusions}\label{sec:conclusions}

We implemented the two released NVIDIA Ising surface code pre-decoders as one integer HLS kernel on four Alveo U55C
cards, each card a node with two kernel instances. The kernel is bit exact with an integer model of the networks on
every instance of every card, and the integer model is evaluated with NVIDIA's unmodified evaluation code. With that
code we reproduce the published accuracy of both networks: the Accurate network lowers the logical error rate of
PyMatching by a factor of \NumRatioAccFpThirteenHigh{} at $p=0.006$, $d=13$, while the Fast network gives little or no
gain at $p=0.003$. Four cards process \NumFourFast{} windows per second of the Fast network at $d=9$ and \NumFourAcc{}
of the Accurate network at $d=13$, with linear scaling over the cards and \NumModelEffFast{}\% and \NumModelEffAcc{}\%
of the throughput predicted by the kernel schedule, but only \NumSpeedFast{} and \NumSpeedAcc{} times the throughput
of the networks in PyTorch on the \NumCpuThreads{} thread host. The latency of a single window, \NumLatFast{}~ms and
\NumLatAcc{}~ms, and the arithmetic the networks need per window show that real time pre-decoding at microsecond
syndrome cycles requires networks designed for hardware rather than faster engines for the released ones. Smaller and
incremental networks, larger engines, resident weights and quantization aware training are the next steps.
